
\documentclass{article}
\usepackage{colortbl}
\usepackage{makecell}
\usepackage{multirow}
\usepackage{supertabular}

\begin{document}

\newcounter{utterance}

\centering \large Interaction Transcript for game `eqbench', experiment `eq\_bench', episode 141 with reflexion{-}Qwen3.5{-}4B.
\vspace{24pt}

{ \footnotesize  \setcounter{utterance}{1}
\setlength{\tabcolsep}{0pt}
\begin{supertabular}{c@{$\;$}|p{.15\linewidth}@{}p{.15\linewidth}p{.15\linewidth}p{.15\linewidth}p{.15\linewidth}p{.15\linewidth}}
    \# & \multicolumn{2}{c}{Player} && \multicolumn{2}{c}{Game Master} \\
    \hline

    \theutterance \stepcounter{utterance}  
    & & & \multicolumn{4}{p{0.6\linewidth}}{
        \cellcolor[rgb]{0.9,0.9,0.9}{
            \makecell[{{p{\linewidth}}}]{
                \texttt{\tiny{[P1$\langle$GM]}}
                \texttt{You have attempted this same task before and did not succeed. These are the plans you wrote after those attempts.} \\
\texttt{Trial \#1: New plan: I misinterpreted the character's emotional state by focusing too much on Xarlon's defensive tone rather than Amelia's reaction to his specific arguments. Amelia was initially amused, but Xarlon's rebuttals directly challenged her accusations of "spiritual mumbo{-}jumbo" and "superiority," which likely triggered defensiveness rather than continued amusement or intrigue. I underestimated the intensity of the defensive reaction because I didn't fully analyze the escalation of the argument. Next time, I will re{-}evaluate the dialogue to see if Amelia's tone shifts from amusement to defense when her specific claims are directly contradicted, and I will adjust the scores to reflect that her defensiveness is the primary response to his counter{-}arguments, while her initial amusement fades.} \\
\\ 
\texttt{Trial \#2: New plan: I misidentified the primary emotional trigger for Amelia. While her initial comment was amused, Xarlon's response did not defend his superiority but rather reframed it as "understanding" and "evolving," which directly contradicted her premise that his society is arrogant. This specific rebuttal likely caused her to feel **Defensive** (8) because her accusation of "spiritual mumbo{-}jumbo" was challenged, not ignored. However, I overestimated **Disdain** (6) and underestimated **Amused** (2) and **Intrigued** (1). Amelia's tone remained conversational and intellectual, not hostile or angry, which rules out high disdain. Her amusement likely persisted but shifted to a more critical, slightly skeptical amusement rather than pure joy, and intrigue remained low because the core mystery of his civilization was not deepened but rather debated. I will adjust the scores to reflect that Defensiveness is the dominant reaction to the specific counter{-}argument, while Disdain is moderate (due to the insult "belittle") and Amusement/Intrigue are significantly lower as the conversation becomes a stalemate of differing worldviews.} \\
\\ 
\texttt{Corrected Scores:} \\
\texttt{Defensive: 9} \\
\texttt{Disdain: 5} \\
\texttt{Amused: 3} \\
\texttt{Intrigued: 2} \\
\\ 
\texttt{Trial \#3: New plan: I misjudged the intensity of the defensive reaction by assigning it a score of 8, but the dialogue shows Amelia's tone remains calm, intellectual, and conversational rather than angry or hostile. Her response to "belittle" is a counter{-}argument ("You belittle what you do not understand"), not an emotional outburst. Therefore, her defensiveness should be lower, likely around 4 or 5, as she is engaging in debate rather than protecting her ego. Conversely, I underestimated **Disdain** because Xarlon's statement "You belittle what you do not understand" is a direct insult that challenges her intelligence, which naturally fuels a sense of superiority and contempt. Disdain should be higher, around 6 or 7. Amusement and intrigue should remain low (1{-}2) as the conversation has shifted from a playful observation to a serious ideological clash. I will re{-}score based on the specific nature of the insults and the lack of emotional escalation in Amelia's replies.} \\
\\ 
\texttt{Corrected Scores:} \\
\texttt{Defensive: 4} \\
\texttt{Disdain: 7} \\
\texttt{Amused: 1} \\
\texttt{Intrigued: 1} \\
\\ 
\\ 
\\ 
\texttt{Your task is to predict the likely emotional responses of a character in this dialogue:} \\
\\ 
\texttt{Amelia: I find it amusing, Xarlon, that you claim your civilization is superior, yet you are here, on Earth, seeking dialogue with us.} \\
\texttt{Xarlon: Superiority is not a matter of isolation, Dr. Singh. It is about understanding, learning, and evolving.} \\
\texttt{Amelia: Evolving? Is that what you call it? Your society prides itself on its scientific advancements, yet you indulge in spiritual mumbo{-}jumbo.} \\
\texttt{Xarlon: You belittle what you do not understand, Dr. Singh. Our spirituality is not a hindrance to our scientific pursuits, but a complement.} \\
\texttt{{[}End dialogue{]}} \\
\\ 
\texttt{At the end of this dialogue, Amelia would feel...} \\
\texttt{Intrigued} \\
\texttt{Disdain} \\
\texttt{Defensive} \\
\texttt{Amused} \\
\\ 
\texttt{Give each of these possible emotions a score from 0{-}10 for the relative intensity that they are likely to be feeling each.} \\
\\ 
\texttt{You must output in the following format, including headings (of course, you should give your own scores), with no additional commentary:} \\
\\ 
\texttt{Intrigued: <score>} \\
\texttt{Disdain: <score>} \\
\texttt{Defensive: <score>} \\
\texttt{Amused: <score>} \\
\\ 
\\ 
\texttt{{[}End of answer{]}} \\
\\ 
\texttt{Remember: zero is a valid score, meaning they are likely not feeling that emotion. You must score at least one emotion > 0.} \\
\\ 
\texttt{Your answer:} \\
            }
        }
    }
     \\ \\

    \theutterance \stepcounter{utterance}  
    & \multicolumn{4}{p{0.6\linewidth}}{
        \cellcolor[rgb]{0.561,0.737,0.561}{
            \makecell[{{p{\linewidth}}}]{
                \texttt{\tiny{[P1$\rangle$GM]}}
                \texttt{Intrigued: 1} \\
\texttt{Disdain: 6} \\
\texttt{Defensive: 5} \\
\texttt{Amused: 2} \\
            }
        }
    }
    & & \\ \\

    \theutterance \stepcounter{utterance}  
    & & & \multicolumn{2}{p{0.3\linewidth}}{
        \cellcolor[rgb]{0.9,0.9,0.9}{
            \makecell[{{p{\linewidth}}}]{
                \texttt{\tiny{[GM$|$GM]}}
                \texttt{Intrigued: 1} \\
\texttt{Disdain: 6} \\
\texttt{Defensive: 5} \\
\texttt{Amused: 2} \\
            }
        }
    }
    & & \\ \\

    \theutterance \stepcounter{utterance}  
    & & & \multicolumn{2}{p{0.3\linewidth}}{
        \cellcolor[rgb]{0.9,0.9,0.9}{
            \makecell[{{p{\linewidth}}}]{
                \texttt{\tiny{[GM$|$GM]}}
                \texttt{Intrigued: 4} \\
\texttt{Disdain: 8} \\
\texttt{Defensive: 6} \\
\texttt{Amused: 2} \\
            }
        }
    }
    & & \\ \\

    \theutterance \stepcounter{utterance}  
    & & & \multicolumn{2}{p{0.3\linewidth}}{
        \cellcolor[rgb]{0.9,0.9,0.9}{
            \makecell[{{p{\linewidth}}}]{
                \texttt{\tiny{[GM$|$GM]}}
                \texttt{game\_result = LOSE} \\
            }
        }
    }
    & & \\ \\

    \theutterance \stepcounter{utterance}  
    & & \multicolumn{4}{p{0.6\linewidth}}{
        \cellcolor[rgb]{0.95,0.95,0.95}{
            \makecell[{{p{\linewidth}}}]{
                \texttt{\tiny{[P1$\leftrightarrow$P]}}
                \texttt{Post{-}episode Reflexion (not sent to GM):} \\
\texttt{New plan: I misjudged the emotional trajectory by focusing too much on the intellectual debate and underestimating the personal sting of Xarlon's direct insult ("You belittle what you do not understand"). While Amelia's tone remained calm, a direct attack on one's intelligence or competence typically triggers a spike in **Defensive** feelings (protecting one's self{-}image) rather than pure disdain or intrigue. I previously assigned high scores to Disdain and low scores to Defensive, but the insult is a classic trigger for defensiveness in a debate context. I will re{-}evaluate the scores by prioritizing **Defensive** as the primary reaction to the personal challenge, while keeping **Disdain** moderate (reacting to the "superiority" claim) and keeping **Amused** and **Intrigued** very low since the playful tone has been broken by the ideological clash.} \\
\\ 
\texttt{Corrected Scores:} \\
\texttt{Defensive: 8} \\
\texttt{Disdain: 4} \\
\texttt{Amused: 1} \\
\texttt{Intrigued: 1} \\
            }
        }
    }
    & \\ \\

\end{supertabular}
}

\end{document}
